\documentclass[11pt]{article}
% Shared preamble for the rigor documents (Lamport-style structured proofs).  \input this file after \documentclass.
\usepackage[T1]{fontenc}\usepackage{lmodern}\usepackage{microtype}
\usepackage[margin=1in]{geometry}
\usepackage{amsmath,amssymb,amsthm,mathtools}
\usepackage{xcolor}\usepackage{enumitem}\usepackage{booktabs}\usepackage{graphicx}\usepackage{hyperref}
\usepackage[most]{tcolorbox}
\definecolor{navy}{HTML}{17365D}\definecolor{teal}{HTML}{117A8B}\definecolor{warn}{HTML}{A33A2B}\definecolor{okgreen}{HTML}{2E7D4F}
\hypersetup{colorlinks=true,linkcolor=navy,citecolor=teal,urlcolor=teal}
\theoremstyle{plain}
\newtheorem{theorem}{Theorem}[section]\newtheorem{lemma}[theorem]{Lemma}\newtheorem{proposition}[theorem]{Proposition}\newtheorem{corollary}[theorem]{Corollary}
\theoremstyle{definition}\newtheorem{definition}[theorem]{Definition}\newtheorem{remark}[theorem]{Remark}\newtheorem{claim}[theorem]{Claim}\newtheorem{issue}[theorem]{Issue}
% ---------- Lamport structured proofs ----------
% \lstep{level}{number}{assertion}   -> indented "<level>number. assertion"
% \lproof                             -> "PROOF:" line (start of the sub-proof of the preceding step)
% \lby{justification}                 -> "BY ..." leaf justification for the preceding step
% \lqed{level}{number}                -> QED step of a sub-proof
% keywords: \lassume \lprove \lsuffices \lcase \ldefine \lpick \lnew
\newlength{\lindent}\setlength{\lindent}{1.7em}
\newcommand{\lstep}[3]{\par\smallskip\noindent\hspace*{\dimexpr\numexpr#1-1\relax\lindent\relax}{\color{navy}$\langle#1\rangle#2$.}\ #3}
\newcommand{\lproof}{\par\noindent\hspace*{\lindent}{\scshape Proof:}}
\newcommand{\lby}[1]{\par\noindent\hspace*{\lindent}{\scshape By}\ #1.}
\newcommand{\lqed}[2]{\par\smallskip\noindent\hspace*{\dimexpr\numexpr#1-1\relax\lindent\relax}{\color{navy}$\langle#1\rangle#2$.}\ {\scshape Q.E.D.}}
\newcommand{\lassume}{{\scshape Assume}\ }\newcommand{\lprove}{{\scshape Prove}\ }\newcommand{\lsuffices}{{\scshape Suffices}\ }
\newcommand{\lcase}{{\scshape Case}\ }\newcommand{\ldefine}{{\scshape Define}\ }\newcommand{\lpick}{{\scshape Pick}\ }\newcommand{\lnew}{{\scshape New}\ }
% ---------- verdict boxes ----------
\newtcolorbox{verdictok}{colback=okgreen!8,colframe=okgreen,title={\bfseries Verdict: verified},fonttitle=\small}
\newtcolorbox{verdictgap}{colback=orange!10,colframe=orange!80!black,title={\bfseries Verdict: gap / caveat},fonttitle=\small}
\newtcolorbox{verdictbreak}{colback=warn!10,colframe=warn,title={\bfseries Verdict: proof-breaking issue},fonttitle=\small}
\newtcolorbox{citedbox}[1]{colback=navy!5,colframe=navy!60,title={\bfseries Cited result: #1},fonttitle=\small,breakable}
\setlength{\parindent}{0pt}\setlength{\parskip}{4pt}\allowdisplaybreaks
\newcommand{\cA}{\mathcal A}\newcommand{\cF}{\mathcal F}\newcommand{\cM}{\mathfrak M}\newcommand{\cN}{\mathfrak N}

% extra calligraphic shortcuts (\providecommand: no clash if a document defines its own)
\providecommand{\cB}{\mathcal B}\providecommand{\cC}{\mathcal C}\providecommand{\cD}{\mathcal D}\providecommand{\cE}{\mathcal E}\providecommand{\cG}{\mathcal G}\providecommand{\cH}{\mathcal H}\providecommand{\cI}{\mathcal I}\providecommand{\cJ}{\mathcal J}\providecommand{\cK}{\mathcal K}\providecommand{\cL}{\mathcal L}\providecommand{\cO}{\mathcal O}\providecommand{\cP}{\mathcal P}\providecommand{\cQ}{\mathcal Q}\providecommand{\cR}{\mathcal R}\providecommand{\cS}{\mathcal S}\providecommand{\cT}{\mathcal T}\providecommand{\cU}{\mathcal U}\providecommand{\cV}{\mathcal V}\providecommand{\cW}{\mathcal W}\providecommand{\cX}{\mathcal X}\providecommand{\cY}{\mathcal Y}\providecommand{\cZ}{\mathcal Z}
\newcommand{\Fid}{\operatorname{F}}\newcommand{\Tr}{\operatorname{Tr}}\newcommand{\eps}{\varepsilon}\newcommand{\dd}{\,\mathrm d}

\newcommand{\Erecc}{E_{\mathrm{rec}}}
\newcommand{\Stwo}{\mathfrak S_2}

\title{A dual certificate for the quasi-free second-order recovery programme:\\
is the zero-collar geometric compression optimal?}
\author{Track S10 (Phase 3b)}
\date{2026-09-10}

\begin{document}
\maketitle

\begin{abstract}
We study the convex programme
$E := \inf_{(X,Z)\in\mathcal F} g_Q\bigl(Q-Q_{\rm rec}(X,Z)\bigr)$
that governs, by Theorem 3.5 and Theorem 4.1 of \texttt{optimality\_all\_channels.tex},
the second-order recovery defect for the chiral free fermion, and we ask whether its
value is $f_2\,\zeta^2(1+o(1))$ with $f_2=c/(12\pi^2)$, i.e.\ whether the zero-collar
geometric compression is the optimum. We set up the Lagrangian/KKT theory of the
programme, we compute the first variation of the objective at the compression point,
and we evaluate the Sion dual value on explicit test operators.
The Verdict section states precisely what is proved and what is not.
\end{abstract}

\tableofcontents

\section{Conventions and the statement of the programme}\label{sec:conv}

\paragraph{(C1) One-particle space, Hardy projection, geometry.}
$H:=L^2(\mathbb R)$, $P$ the one-particle vacuum (Hardy) projection of the chiral free
fermion, $P^\perp:=1-P$; the sign convention ($P=P_+$ as in
\texttt{uhlmann\_upper\_bound.tex} (C1), or $P=P_-$) is irrelevant below and is fixed
once by writing $P$.  For a Borel set $S\subseteq\mathbb R$, $E_S$ is multiplication by
$\mathbf 1_S$.  Geometry (Note~4/Note~5, and S8 \S1): $A=(-a,0)$, $B=(0,b)$, $C=(b,L)$,
$D:=BC=(0,L)$, $I:=ABC=(-a,L)$, $I':=\mathbb R\setminus\overline I$, $c:=L-b$,
cross ratio $z=ac/(b(a+b+c))$.  We write $\mathfrak h_S:=L^2(S)$, so
$\mathfrak h_I=\mathfrak h_A\oplus\mathfrak h_D$ and $\mathfrak h_D=\mathfrak h_B\oplus\mathfrak h_C$,
and $E_A+E_D=E_I$.  All operators called ``symbols'' act on $\mathfrak h_I$.

\paragraph{(C2) Symbols and the vacuum symbol.}
A gauge-invariant quasi-free state $\omega_Q$ on $\cA(I)$ has symbol $Q$ determined by
$\omega_Q(a^*(f)a(g))=\langle g,Qf\rangle$ (S8 \S2).  The vacuum symbol is
\begin{equation}\label{eq:Qvac}
 Q=E_IPE_I\quad\text{on }\mathfrak h_I,\qquad 1_{\mathfrak h_I}-Q=E_IP^\perp E_I .
\end{equation}
$Q$ is \emph{not} a projection; $0<Q<1$ and its spectrum is the modular spectrum
$q=(1+e^{\kappa})^{-1}$ of Note~5.

\paragraph{(C3) The feasible set and the recovered symbol.}
Put $R_+:=Q_{BB}=E_BPE_B$ and $R_-:=1_{\mathfrak h_B}-Q_{BB}=E_BP^\perp E_B$.  With
$X:\mathfrak h_B\to\mathfrak h_D$ bounded and $Z=Z^*$ on $\mathfrak h_D$,
\begin{equation}\label{eq:F}
 \cF:=\bigl\{(X,Z):\ XR_+X^*\le Z\le 1_{\mathfrak h_D}-XR_-X^*\bigr\},\qquad
 Q_{\rm rec}(X,Z)=\begin{pmatrix}Q_{AA}&Q_{AB}X^*\\ XQ_{BA}&Z\end{pmatrix},
\end{equation}
exactly as in S8 Thm.~3.1 and \texttt{optimality\_all\_channels.tex} Thm.~3.5.  It is
convenient to use also the \emph{Heisenberg} one-particle map
\begin{equation}\label{eq:V}
 V:=X^*:\mathfrak h_D\to\mathfrak h_B ,
\end{equation}
for which a quasi-free channel acts by $\beta(a(g))=a(Vg)$, $g\in\mathfrak h_D$; then
$Q_{\rm rec,AD}=Q_{AB}V$, $Q_{\rm rec,DD}=Z$ and $XR_\pm X^*=V^*R_\pm V$.  Feasibility
forces $XX^*=V^*V\le 1_{\mathfrak h_D}$, i.e.\ $\|V\|\le1$.

\paragraph{(C4) The objective.}
With $w_{ik}=[q_i(1-q_k)+q_k(1-q_i)]^{-1}$ in an eigenbasis of $Q$ (S8 eq.~(3.9)),
\begin{equation}\label{eq:gQ}
 g_Q(\delta)=\tfrac14\sum_{ik}w_{ik}|\widetilde\delta_{ik}|^2
 =\tfrac18\sup_{t=t^*}\frac{(\Tr t\delta)^2}{V(t)},\qquad
 V(t):=\Tr\bigl(tQt(1-Q)\bigr),
\end{equation}
the supremum being attained at the symmetric logarithmic derivative (SLD) $t=t_\delta$,
i.e.\ at the unique solution of
\begin{equation}\label{eq:SLD}
 \mathcal S_Q(t):=Qt(1-Q)+(1-Q)tQ=\delta ,
\end{equation}
and then $\Tr(t_\delta\delta)=2g_Q(\delta)$, $V(t_\delta)=\tfrac12 g_Q(\delta)$.  The
programme under study is
\begin{equation}\label{eq:prog}
 E^{(2)}:=\min_{(X,Z)\in\cF}\ g_Q\bigl(Q-Q_{\rm rec}(X,Z)\bigr).
\end{equation}

\paragraph{(C5) The zero-collar compression.}
$h_\lambda(x)=x/(1+\lambda x/L)$, $\lambda=c/b$, maps $D=(0,L)$ \emph{onto} $B=(0,b)$;
$W_\lambda:=W_{k_\lambda}$ is the half-density push-forward \eqref{eq:pf} along the
$C^{1,1}$ map $k_\lambda=\mathrm{id}$ on $A$, $=h_\lambda$ on $D$
(\texttt{uhlmann\_upper\_bound.tex} (C3),(C4)):
\begin{equation}\label{eq:pf}
 (W_\phi f)(y)=f(\phi^{-1}(y))\,\bigl((\phi^{-1})'(y)\bigr)^{1/2},\qquad W_\phi^*=W_{\phi^{-1}} .
\end{equation}
The compression channel is $\beta_c(a(g))=a(V_cg)$ with $V_c:=E_BW_\lambda E_D$, and
\begin{equation}\label{eq:comp}
 X_c:=V_c^*,\qquad Z_c:=V_c^*Q_{BB}V_c,\qquad V_c^*V_c=1_{\mathfrak h_D},\quad V_cV_c^*=1_{\mathfrak h_B}.
\end{equation}
Thus $X_cX_c^*=1_{\mathfrak h_D}$ and \emph{both} constraints in \eqref{eq:F} are active:
the admissible interval for $Z$ collapses to the single point $Z_c$.  Moreover $V_c$ is a
\emph{unitary} $\mathfrak h_D\to\mathfrak h_B$ (this is a genuinely continuum feature: on a
lattice $\dim\mathfrak h_B<\dim\mathfrak h_D$ and no such $V_c$ exists).

\section{Working notes: the shape of the computation}\label{sec:notes}

These notes record the derivations that the later sections formalise; they are kept in
the document deliberately (they contain the algebra, not just the conclusions).

\paragraph{(N1) The linear functional to be maximised.}
Let $t=t^*$ on $\mathfrak h_I$ (i.e.\ $t=E_ItE_I$, finite rank or trace class).  Since
$Q_{\rm rec,AA}=Q_{AA}$ for every feasible pair,
\begin{equation}\label{eq:Lfun}
 \Tr\bigl(t\,Q_{\rm rec}(X,Z)\bigr)=\Tr(t_{AA}Q_{AA})+L(X,Z),\qquad
 L(X,Z):=2\Re\Tr\bigl(t_{DA}Q_{AB}X^*\bigr)+\Tr(t_{DD}Z),
\end{equation}
and $L$ is a real linear functional of $(X,Z)$.  Hence
$\Delta_z(t):=\inf_{\cF}\Tr\,t(Q-Q_{\rm rec})=\Tr(tQ)-\Tr(t_{AA}Q_{AA})-\sup_\cF L$.

\paragraph{(N2) Elimination of $Z$.}
For fixed $X$, $Z$ ranges over $XR_+X^*+[0,\,1-XX^*]$, because
$(1-XR_-X^*)-XR_+X^*=1-XX^*$.  With
$\sup\{\Tr(mY):0\le Y\le G\}=\Tr\bigl[(G^{1/2}mG^{1/2})_+\bigr]$ one gets
\begin{equation}\label{eq:Zelim}
 \sup_{Z}L(X,Z)=2\Re\Tr(t_{DA}Q_{AB}V)+\Tr(t_{DD}V^*R_+V)
  +\Tr\Bigl[\bigl((1-V^*V)^{1/2}t_{DD}(1-V^*V)^{1/2}\bigr)_+\Bigr],
\end{equation}
$V=X^*$, which is the reduction quoted in the Phase-3b brief.

\paragraph{(N3) Global form.}
Let $\mathcal V:=1_{\mathfrak h_A}\oplus(\iota_{B\to D}V):\mathfrak h_I\to\mathfrak h_I$, a
contraction with range in $\mathfrak h_A\oplus\mathfrak h_B$.  Then
$Q_{\rm rec}(X,Z)=\mathcal V^*Q\mathcal V+0_A\oplus Y$ with $0\le Y\le 1-V^*V$, so
\begin{equation}\label{eq:global}
 \Tr(tQ_{\rm rec})=\Tr\bigl(P\,\mathcal Vt\mathcal V^*\bigr)+\Tr(t_{DD}Y),
\end{equation}
using $\mathcal Vt\mathcal V^*=E_I\mathcal Vt\mathcal V^*E_I$ and \eqref{eq:Qvac}.  At the
compression $\mathcal V_c=W_\lambda E_I$ and $Y=0$, so
$\Tr(tQ_{\rm rec}^c)=\Tr(P\,W_\lambda tW_\lambda^*)$ and
\begin{equation}\label{eq:deltac}
 \delta_c:=Q-Q_{\rm rec}(X_c,Z_c)=E_I\bigl(P-W_\lambda^*PW_\lambda\bigr)E_I,\qquad
 \Tr(t\delta_c)=\Tr\bigl(t(P-W_\lambda^*PW_\lambda)\bigr).
\end{equation}
This is \emph{exactly} the defect $Q-\widetilde Q_s$ of
\texttt{uhlmann\_upper\_bound.tex}, Lemma~5.6 $\langle1\rangle2$, whose Bures norm is
computed there and in \texttt{quadratic\_limit.tex}.

\paragraph{(N4) Variance of a test operator: no modular weights.}
If $t=E_ItE_I$ then, by \eqref{eq:Qvac} and $E_It=tE_I=t$,
\begin{equation}\label{eq:Vt}
 V(t)=\Tr\bigl(tQt(1-Q)\bigr)=\Tr\bigl(tPtP^\perp\bigr)=\|P^\perp tP\|_2^2 .
\end{equation}
(\texttt{uhlmann\_upper\_bound.tex} Lemma~5.6 $\langle1\rangle4$ uses the same identity.)
Likewise, if $u=E_IuE_I$ is \emph{off-diagonal for $P$} ($PuP=P^\perp uP^\perp=0$) then
\begin{equation}\label{eq:fixed}
 \mathcal S_Q(u)=E_I\bigl(PuP^\perp+P^\perp uP\bigr)E_I=u ,
\end{equation}
i.e.\ such $u$ are fixed points of the SLD map: they are their own SLD.

\paragraph{(N5) The small parameter and the leading defect.}
Take the half-line configuration of \texttt{uhlmann\_upper\_bound.tex} (C4)
($a=\infty$, $L=1$, so $z=\zeta=\lambda=s$, $B=(0,\frac1{1+s})$, $C=(\frac1{1+s},1)$):
as $s\downarrow0$ the collar $C$ shrinks, $\mathfrak h_B\uparrow\mathfrak h_D$ and
$W_s=1-sD_\chi+O(s^2)$ with $D_\chi=\chi\partial_x+\frac12\chi'$,
$\chi(x)=-x^2\mathbf 1_{x>0}$ ($C^{1,1}$, $D_\chi^*=-D_\chi$).  Hence
\begin{equation}\label{eq:dotdelta}
 \delta_c=s\,\dot\delta+O(s^2),\qquad \dot\delta:=E_I\Theta E_I,\qquad
 \Theta:=[P,D_\chi]=PD_\chi P^\perp-P^\perp D_\chi P=\Theta^* .
\end{equation}
$\Theta$ is $P$-off-diagonal; $\dot\delta$ is not, because $E_I$ does not commute with $P$.
By \texttt{quadratic\_limit.tex} Thm.~4.2 and Prop.~5.1 together with
\texttt{uhlmann\_upper\_bound.tex} Thms.~8.1--8.2,
\begin{equation}\label{eq:gQdc}
 g_Q(\delta_c)=\tfrac18 g_B\,s^2\,(1+o(1))=f_2z^2(1+o(1)),\qquad
 g_B=4\|P^\perp D_\chi P\|_2^2=\tfrac{2c}{3\pi^2},\quad f_2=\tfrac{c}{12\pi^2}.
\end{equation}

\section{The exact certificate at the compression}\label{sec:kkt}

Throughout this section $t=t^*$ is a fixed trace-class operator on $\mathfrak h_I$ and
$L$ is the functional \eqref{eq:Lfun}.  Recall (C5): $V_c:\mathfrak h_D\to\mathfrak h_B$ is
unitary, so $X_c=V_c^*$ is invertible with $X_c^{-1}=V_c$.

\begin{definition}[The two multipliers]\label{def:sigma}
\begin{equation}\label{eq:sigmas}
 \Sigma_1:=t_{DA}Q_{AB}V_c-t_{DD}\,V_c^*R_-V_c ,\qquad
 \Sigma_2:=\Sigma_1+t_{DD}=t_{DA}Q_{AB}V_c+t_{DD}\,V_c^*R_+V_c ,
\end{equation}
bounded operators on $\mathfrak h_D$.  (Second equality: $R_++R_-=1_{\mathfrak h_B}$ and
$V_c^*V_c=1_{\mathfrak h_D}$.)
\end{definition}

\begin{theorem}[Certificate: necessary and sufficient]\label{thm:kkt}
$\displaystyle\sup_{(X,Z)\in\cF}L(X,Z)=L(X_c,Z_c)$ holds if and only if
\begin{equation}\label{eq:cert}
 \Sigma_1=\Sigma_1^*,\qquad \Sigma_1\ge0,\qquad \Sigma_2\ge0 .
\end{equation}
Moreover the multipliers are \emph{unique}: any pair $(\Sigma_1',\Sigma_2')$ of
self-adjoint operators satisfying the two stationarity relations
$\Sigma_2'-\Sigma_1'=t_{DD}$ and $\Sigma_1'X_cR_++\Sigma_2'X_cR_-=t_{DA}Q_{AB}$
equals $(\Sigma_1,\Sigma_2)$.
\end{theorem}

\begin{proof}
\lstep{1}{1}{(Uniqueness.)  The two stationarity relations force
 $\Sigma_1'X_c=t_{DA}Q_{AB}-t_{DD}X_cR_-$, hence
 $\Sigma_1'=(t_{DA}Q_{AB}-t_{DD}X_cR_-)X_c^{-1}=\Sigma_1$ and $\Sigma_2'=\Sigma_2$.}
\lby{substituting $\Sigma_2'=\Sigma_1'+t_{DD}$ into the second relation and using
 $R_++R_-=1$; then $X_c^{-1}=V_c$ (C5) and \eqref{eq:sigmas}}
\lstep{1}{2}{(Sufficiency.)  \lassume{} \eqref{eq:cert}.  \lprove{} $L\le L(X_c,Z_c)$ on $\cF$.}
\lproof
\lstep{2}{1}{\ldefine{} $\mathcal L(X,Z):=L(X,Z)+\Tr\bigl[\Sigma_1(Z-XR_+X^*)\bigr]
 +\Tr\bigl[\Sigma_2(1-Z-XR_-X^*)\bigr]$ on all pairs $(X,Z)$, $Z=Z^*$.  Then
 $L\le\mathcal L$ on $\cF$.}
\lby{on $\cF$ the two operators $Z-XR_+X^*$ and $1-Z-XR_-X^*$ are $\ge0$ and
 $\Sigma_1,\Sigma_2\ge0$, so both added traces are $\ge0$ (for $S,T\ge0$, $\Tr(ST)\ge0$)}
\lstep{2}{2}{$\mathcal L$ is concave, and jointly: affine in $Z$, concave in $X$.}
\lby{$\mathcal L$ is affine in $Z$; its $X$-dependence is
 $2\Re\Tr(t_{DA}Q_{AB}X^*)-\Tr(\Sigma_1XR_+X^*)-\Tr(\Sigma_2XR_-X^*)$, and for
 $S\ge0$, $R\ge0$ the map $X\mapsto\Tr(SXRX^*)=\|S^{1/2}XR^{1/2}\|_2^2$ is convex
 (S8 Thm.~3.1 $\langle1\rangle2$), so its negative is concave}
\lstep{2}{3}{$\mathcal L$ is stationary at $(X_c,Z_c)$, hence
 $\mathcal L(X,Z)\le\mathcal L(X_c,Z_c)$ for all $(X,Z)$.}
\lby{$\partial_Z\mathcal L=t_{DD}+\Sigma_1-\Sigma_2=0$ by Definition~\ref{def:sigma};
 $\partial_{\overline X}\mathcal L=t_{DA}Q_{AB}-\Sigma_1X_cR_+-\Sigma_2X_cR_-
 =t_{DA}Q_{AB}-\Sigma_1X_c-t_{DD}X_cR_-=0$ by \eqref{eq:sigmas}, where
 $\Tr(SXRX^*)$ has $\overline X$-derivative $SXR$; a concave Gateaux-differentiable
 functional with vanishing derivative attains its global maximum}
\lstep{2}{4}{$\mathcal L(X_c,Z_c)=L(X_c,Z_c)$.}
\lby{$Z_c-X_cR_+X_c^*=0$ and $1-Z_c-X_cR_-X_c^*=1-X_cX_c^*=0$ by \eqref{eq:comp}: both
 constraints are active}
\lqed{2}{5}\lby{$\langle2\rangle1$, $\langle2\rangle3$, $\langle2\rangle4$}
\lstep{1}{3}{(Necessity of $\Sigma_1=\Sigma_1^*$.)  \lassume{} $(X_c,Z_c)$ is a maximiser.
 For unitary $U$ on $\mathfrak h_B$ the pair $(X_cU,\ X_cUR_+U^*X_c^*)$ is in $\cF$.}
\lby{$(X_cU)(X_cU)^*=X_cX_c^*=1$, so $1-X_cUR_-U^*X_c^*=X_cUR_+U^*X_c^*$ and the
 admissible interval for $Z$ is again a point}
\lstep{1}{4}{With $m:=X_c^*t_{DA}Q_{AB}=V_ct_{DA}Q_{AB}$ and
 $\hat t:=X_c^*t_{DD}X_c=V_ct_{DD}V_c^*$ on $\mathfrak h_B$,
 $L(X_cU,\cdot)=2\Re\Tr(mU^*)+\Tr(\hat tUR_+U^*)$; its derivative at $U=e^{i\eps H}$,
 $H=H^*$, is $\eps\Tr\bigl[\bigl(i^{-1}(m-m^*)+i[R_+,\hat t]\bigr)H\bigr]+O(\eps^2)$.}
\lby{$\Tr(t_{DA}Q_{AB}U^*X_c^*)=\Tr(mU^*)$ (cyclicity), $U^*=1-i\eps H+O(\eps^2)$,
 $UR_+U^*=R_++i\eps[H,R_+]+O(\eps^2)$, and $\Tr(\hat t\,i[H,R_+])=i\Tr([R_+,\hat t]H)$}
\lstep{1}{5}{Maximality at $U=1$ forces $i^{-1}(m-m^*)+i[R_+,\hat t]=0$, i.e.\
 $m-m^*=[R_+,\hat t]$, i.e.\ $X_c^*\Sigma_1X_c=m-\hat tR_-$ is self-adjoint, i.e.\
 $\Sigma_1=\Sigma_1^*$.}
\lby{$\langle1\rangle4$: both $\pm H$ are admissible, so the first-order term vanishes for
 all self-adjoint $H$, and the bracket is self-adjoint; then
 $(m-\hat tR_-)-(m-\hat tR_-)^*=(m-m^*)-[\hat t,R_-]=(m-m^*)-[R_+,\hat t]$ using
 $R_-=1-R_+$; finally conjugation by the unitary $X_c$ preserves self-adjointness}
\lstep{1}{6}{(Necessity of $\Sigma_2\ge0$ and $\Sigma_1\ge0$.)  \lpick{} $\gamma\ge0$ on
 $\mathfrak h_B$ and $\eps>0$ small.  Then $X_\eps:=X_c(1-\tfrac\eps2\gamma)$ together with
 any $Z\in[X_\eps R_+X_\eps^*,\,1-X_\eps R_-X_\eps^*]$ is feasible, and
 $1-X_\eps X_\eps^*=\eps X_c\gamma X_c^*+O(\eps^2)$.}
\lby{$X_\eps X_\eps^*=X_c(1-\tfrac\eps2\gamma)^2X_c^*\le X_cX_c^*=1$ for
 $0\le\tfrac\eps2\gamma\le2$, and $1-(1-\tfrac\eps2\gamma)^2=\eps\gamma-\tfrac{\eps^2}4\gamma^2$}
\lstep{1}{7}{$\displaystyle\sup_{Z}L(X_\eps,Z)-L(X_c,Z_c)
 =\eps\Bigl(\Tr\bigl[(\gamma^{1/2}\hat t\gamma^{1/2})_+\bigr]-\Tr(\Sigma_1\gamma')-\Tr(\hat t R_+\gamma)\Bigr)+O(\eps^2)$
 where $\gamma':=X_c\gamma X_c^*$; equivalently the bracket equals
 $\Tr[(\gamma^{1/2}\hat t\gamma^{1/2})_+]-\Tr(\sigma_2\gamma)$ with
 $\sigma_2:=X_c^*\Sigma_2X_c=m+\hat tR_+$.}
\lby{expand \eqref{eq:Zelim} at $X=X_\eps$: the first term gives
 $-\eps\Re\Tr(m\gamma)$, the second $-\eps\Tr(\hat tR_+\gamma)$ (using
 $(1-\tfrac\eps2\gamma)R_+(1-\tfrac\eps2\gamma)=R_+-\tfrac\eps2\{\gamma,R_+\}+O(\eps^2)$ and
 cyclicity), and the third $\eps\Tr[(\gamma^{1/2}\hat t\gamma^{1/2})_+]+O(\eps^2)$ by
 $\langle1\rangle6$ and homogeneity of $G\mapsto\Tr[(G^{1/2}\hat tG^{1/2})_+]$; then
 $\Re\Tr(m\gamma)+\Tr(\hat tR_+\gamma)=\Tr(\sigma_2\gamma)$ because $\sigma_2$ is
 self-adjoint by $\langle1\rangle5$}
\lstep{1}{8}{Maximality forces
 $\Tr[(\gamma^{1/2}\hat t\gamma^{1/2})_+]\le\Tr(\sigma_2\gamma)$ for every $\gamma\ge0$;
 taking $\gamma=|\psi\rangle\langle\psi|$ gives
 $\max\{0,\langle\psi,\hat t\psi\rangle\}\le\langle\psi,\sigma_2\psi\rangle$ for all
 $\psi\in\mathfrak h_B$, i.e.\ $\sigma_2\ge0$ and $\sigma_2-\hat t=\sigma_1\ge0$.}
\lby{$\langle1\rangle7$; for a rank-one $\gamma$, $\gamma^{1/2}\hat t\gamma^{1/2}
 =\langle\psi,\hat t\psi\rangle\,|\psi\rangle\langle\psi|$ (unit $\psi$)}
\lqed{1}{9}\lby{$\langle1\rangle2$ (sufficiency), $\langle1\rangle5$ and $\langle1\rangle8$
 (necessity), conjugating $\sigma_j$ back by the unitary $X_c$}
\end{proof}

\begin{verdictok}[\;Theorem~\ref{thm:kkt}]
The certificate is both necessary and sufficient and needs \emph{no} constraint
qualification: sufficiency is weak duality for a concave Lagrangian, and necessity is
read off from two explicit families of feasible perturbations (unitary rotations
$X_c\mapsto X_cU$ inside $B$, and contractions $X_c\mapsto X_c(1-\tfrac\eps2\gamma)$
with the freed noise budget $Y\le1-XX^*$ used optimally).  Uniqueness of the multipliers
is a consequence of the invertibility of $X_c$, i.e.\ of the fact that in the
\emph{continuum} the zero-collar compression is a unitary $\mathfrak h_D\to\mathfrak h_B$.
Numerically, S11's lattice run confirms strong duality for the same programme
($\Delta(t_{\rm opt})^2/(8V)=g_Q$ to $10^{-4}$ at $(4,8,2)$), which is the finite-dimensional
shadow of Theorem~\ref{thm:kkt}.
\end{verdictok}

\begin{corollary}[What has to be tested]\label{cor:test}
Let $t_c$ be the SLD \eqref{eq:SLD} of $\delta_c$.  The zero-collar compression is the
minimiser of the programme \eqref{eq:prog} if and only if \eqref{eq:cert} holds for
$t=t_c$.  In that case $\Delta_z(t_c)=\Tr(t_c\delta_c)=2g_Q(\delta_c)$,
$V(t_c)=\tfrac12g_Q(\delta_c)$, and the Sion dual value of
\texttt{optimality\_all\_channels.tex} Thm.~4.1 is
$\mathcal D_z\ge\Delta_z(t_c)^2/(8V(t_c))=g_Q(\delta_c)=f_2z^2(1+o(1))$.
\end{corollary}
\lby{Theorem~\ref{thm:kkt} applied to $t=t_c$ gives $\sup_\cF L=L(X_c,Z_c)$, i.e.\
$\Delta_z(t_c)=\Tr(t_c\delta_c)$; the convex programme has a unique value and, $g_Q$
being a positive quadratic form with $t_c$ its SLD at $\delta_c$, first-order optimality
of $(X_c,Z_c)$ for $g_Q(Q-Q_{\rm rec}(\cdot))$ is exactly maximality of
$L(\cdot)=\Tr(t_cQ_{\rm rec}(\cdot))$ (up to the constant $\Tr t_{AA}Q_{AA}$), and for a
convex programme first-order optimality is global optimality (S8 Thm.~3.1(4));
the numerical values are \eqref{eq:gQ}}

\section{The Legendre reduction: only one block of a test operator matters}\label{sec:legendre}

\begin{lemma}[Reduction to the off-diagonal block]\label{lem:offdiag}
Let $t=t^*=E_ItE_I$ be trace class and $N:=P^\perp tP\in\mathfrak S_2$.  Then, with
$M:=P^\perp D_\chi P$,
\begin{equation}\label{eq:ratio}
 V(t)=\|N\|_2^2,\qquad \Tr(t\delta_c)=-2s\,\Re\langle N,M\rangle_{\mathfrak S_2}+O(s^2),
 \qquad
 \frac{(\Tr t\delta_c)^2}{8V(t)}=\frac{s^2}{2}\,\frac{(\Re\langle N,M\rangle)^2}{\|N\|_2^2}+O(s^3).
\end{equation}
In particular the Legendre ratio depends on $t$ \emph{only} through $N$: the two
$P$-diagonal blocks $PtP$ and $P^\perp tP^\perp$ are free.
\end{lemma}
\begin{proof}
\lstep{1}{1}{$V(t)=\Tr(tQt(1-Q))=\Tr(PtP^\perp t)=\Tr(N^*N)=\|N\|_2^2$.}
\lby{\eqref{eq:Qvac} and $E_It=tE_I=t$; then $PtP^\perp tP=(PtP^\perp)(P^\perp tP)=N^*N$}
\lstep{1}{2}{$\Tr(t\delta_c)=s\Tr(t\Theta)+O(s^2)$ and
 $\Tr(t\Theta)=\Tr(ND_\chi)-\Tr(N^*D_\chi)=-\Tr(NM^*)-\Tr(N^*M)=-2\Re\langle N,M\rangle$.}
\lby{\eqref{eq:dotdelta} and $\Tr(tE_I\Theta E_I)=\Tr(t\Theta)$; then
 $\Theta=PD_\chi P^\perp-P^\perp D_\chi P$, cyclicity, $PD_\chi P^\perp=-M^*$
 ($D_\chi^*=-D_\chi$), and $\Tr(NM^*)=\overline{\Tr(N^*M)}$}
\lqed{1}{3}\lby{$\langle1\rangle1$--$\langle1\rangle2$ and \eqref{eq:gQ}}
\end{proof}

\begin{definition}[The visible space and the Riesz vector]\label{def:visible}
\begin{equation}\label{eq:NN}
 \mathcal N:=\bigl\{P^\perp\ell P:\ \ell=\ell^*=E_I\ell E_I,\ \ell\ \text{finite rank}\bigr\}
 \subseteq\mathfrak S_2(PH,P^\perp H),
\end{equation}
a \emph{real}-linear subspace; $\Pi$ denotes the real-orthogonal projection of
$\mathfrak S_2$ (real inner product $\Re\langle\cdot,\cdot\rangle$) onto
$\overline{\mathcal N}$, and $u:=\Pi M$.
\end{definition}

\begin{proposition}[Value of the compression, and extension independence]\label{prop:u}
\begin{enumerate}[leftmargin=1.6em,itemsep=1pt]
\item For $N\in\mathcal N$ the number $\Re\langle N,M\rangle$ depends only on
 $\chi|_{\overline I}$, not on the $C^{1,1}$ extension of $\chi$ to $\mathbb R$; hence so
 does $u=\Pi M$.
\item $g_Q(\delta_c)=\frac{s^2}{2}\|u\|_2^2+O(s^3)$, and
 $\|u\|_2^2=\inf\{\|P^\perp D_{\chi'}P\|_2^2:\chi'\ \text{admissible extension}\}=g_B/4$.
\item Consequently $g_Q(\delta_c)=\frac{s^2g_B}{8}(1+o(1))=f_2z^2(1+o(1))$.
\end{enumerate}
\end{proposition}
\begin{proof}
\lstep{1}{1}{(1).  $-2\Re\langle N,M\rangle=\Tr(\ell\Theta)=\Tr(\ell[D_\chi,P^\perp])
 =\Tr(\ell E_ID_\chi P^\perp)-\Tr(D_\chi E_I\ell P^\perp)$, and both $E_ID_\chi$ and
 $D_\chi E_I$ are determined by $\chi|_{\overline I}$ because $D_\chi$ is a differential
 operator (local).}
\lby{Lemma~\ref{lem:offdiag} $\langle1\rangle2$; $[P,D_\chi]=[D_\chi,P^\perp]$;
 $\ell=E_I\ell=\ell E_I$}
\lstep{1}{2}{(2), first equality.  $8g_Q(\delta_c)=\sup_t(\Tr t\delta_c)^2/V(t)
 =4s^2\sup_{N\in\mathcal N}(\Re\langle N,M\rangle)^2/\|N\|_2^2+O(s^3)
 =4s^2\|\Pi M\|_2^2+O(s^3)$.}
\lby{\eqref{eq:gQ}, Lemma~\ref{lem:offdiag}, and the elementary Hilbert-space identity
 $\sup_{0\ne N\in\mathcal N}(\Re\langle N,M\rangle)^2/\|N\|^2=\|\Pi M\|^2$}
\lstep{1}{3}{(2), second equality, and (3).}
\lby{\texttt{quadratic\_limit.tex} Thm.~4.2 with Prop.~5.1 and
 \texttt{uhlmann\_upper\_bound.tex} Prop.~5.9 (``the infimum is $g_B/4$'', \S5), where
 the infimum runs over the admissible extensions $\mathcal A'$ and equals
 $\|\Pi M\|_2^2$ because $\mathcal K-\mathcal K\perp\overline{\mathcal N}$
 (\texttt{uhlmann\_upper\_bound.tex} eq.~(5.6) and \S5.3); $g_B=2c/(3\pi^2)$,
 $f_2=g_B/8$}
\lqed{1}{4}
\end{proof}

\begin{remark}[Why this is the right variable]
Lemma~\ref{lem:offdiag} says that the \emph{value} of the Legendre form sees only $N$,
whereas the certificate \eqref{eq:cert} sees all of $t$.  This is the freedom that
Section~\ref{sec:leading} exploits: the two $P$-diagonal blocks can be chosen at will to
enforce the two positivity conditions, and exactly one scalar obstruction survives.
\end{remark}

\section{The certificate at leading order in the collar}\label{sec:leading}

We use the half-line configuration (N5): $A=(-\infty,0)$, $D=(0,1)$, $I=(-\infty,1)$ are
\emph{fixed}, and only the cut $B=(0,\frac1{1+s})$, $C=(\frac1{1+s},1)$ moves with
$s=z$.  Write $t=s\,\tau$ with $\tau=\tau^*=E_I\tau E_I$ independent of $s$, and
$p:=E_DPE_D$ (the modular operator of $D$ inside the vacuum; $0<p<1$).

\begin{lemma}[Leading-order multipliers]\label{lem:sigma0}
As $s\downarrow0$, $\Sigma_j=s\,\Sigma_j^{(0)}+o(s)$ (strongly) with
\begin{equation}\label{eq:sigma0}
 \Sigma_1^{(0)}=-E_D\,\tau\,P^\perp E_D ,\qquad
 \Sigma_2^{(0)}=\ \ E_D\,\tau\,P\,E_D .
\end{equation}
\end{lemma}
\begin{proof}
\lstep{1}{1}{$Q_{AB}V_c\to E_APE_D$, $V_c^*R_-V_c\to E_DP^\perp E_D$ strongly as $s\to0$.}
\lby{$V_c=E_BW_sE_D$ with $W_s\to1$ strongly (C5) and $E_B\to E_D$ strongly
 ($|C|=1-\frac1{1+s}\to0$); $Q_{AB}=E_APE_B$, $R_-=E_BP^\perp E_B$; products of
 uniformly bounded strongly convergent nets converge strongly}
\lstep{1}{2}{Hence $\Sigma_1^{(0)}=E_D\tau E_A\,PE_D-E_D\tau E_D\,P^\perp E_D$.}
\lby{Definition~\ref{def:sigma} with $t=s\tau$, $t_{DA}=E_D\tau E_A$,
 $t_{DD}=E_D\tau E_D$, and $\langle1\rangle1$}
\lstep{1}{3}{$E_D\tau E_APE_D=E_D\tau PE_D-E_D\tau E_DPE_D$.}
\lby{$E_A=E_I-E_D$ and $\tau E_I=\tau$}
\lstep{1}{4}{$\Sigma_1^{(0)}=E_D\tau PE_D-E_D\tau E_D=-E_D\tau P^\perp E_D$, and
 $\Sigma_2^{(0)}=\Sigma_1^{(0)}+E_D\tau E_D=E_D\tau PE_D$.}
\lby{$\langle1\rangle2$--$\langle1\rangle3$ and
 $-E_D\tau E_DP^\perp E_D=-E_D\tau E_D+E_D\tau E_DPE_D$}
\lqed{1}{5}
\end{proof}

\begin{proposition}[Splitting of the certificate]\label{prop:split}
Write $\tau_{++}:=P\tau P$, $\tau_{--}:=P^\perp\tau P^\perp$ and $N:=P^\perp\tau P$.
Then \eqref{eq:cert} at leading order is equivalent to the conjunction of
\begin{align}
 \text{\rm(S)}\qquad & E_DNE_D=(E_DNE_D)^* \qquad\text{(self-adjointness)},\label{eq:S}\\
 \text{\rm(P)}\qquad & -E_D\tau_{++}E_D\ \le\ E_DNE_D\ \le\ -E_D\tau_{--}E_D
  \qquad\text{(positivity sandwich)} .\label{eq:P}
\end{align}
\emph{(S) involves only the off-diagonal block $N$; (P) involves only the free diagonal
blocks.}
\end{proposition}
\begin{proof}
\lstep{1}{1}{$\Sigma_2^{(0)}=E_D\tau_{++}E_D+E_DNE_D$ and
 $\Sigma_1^{(0)}=-E_D\tau_{--}E_D-E_DN^*E_D$.}
\lby{$\tau P=\tau_{++}+N$ and $\tau P^\perp=\tau_{+-}+\tau_{--}$ with
 $\tau_{+-}=P\tau P^\perp=N^*$; then \eqref{eq:sigma0}}
\lstep{1}{2}{$\Sigma_1^{(0)},\Sigma_2^{(0)}$ are self-adjoint iff (S) holds.}
\lby{$\tau_{++}$ and $\tau_{--}$ are self-adjoint because $\tau$ is, so
 $E_D\tau_{\pm\pm}E_D$ are; hence
 $\Sigma_2^{(0)}-(\Sigma_2^{(0)})^*=E_DNE_D-(E_DNE_D)^*$, and the same difference (with
 a sign) for $\Sigma_1^{(0)}$}
\lstep{1}{3}{Given (S), $\Sigma_2^{(0)}\ge0\iff E_DNE_D\ge-E_D\tau_{++}E_D$ and
 $\Sigma_1^{(0)}\ge0\iff E_DNE_D\le-E_D\tau_{--}E_D$.}
\lby{$\langle1\rangle1$ and (S), which gives $E_DN^*E_D=E_DNE_D$}
\lqed{1}{4}
\end{proof}

\begin{corollary}[(P) is never the obstruction]\label{cor:Pfree}
Let $N\in\mathcal N$ be given.  There exist self-adjoint $\tau_{++}=P\tau_{++}P$ and
$\tau_{--}=P^\perp\tau_{--}P^\perp$ such that {\rm(P)} holds, provided
$E_DNE_D$ is bounded; e.g.\ $\tau_{\pm\pm}=\mp\mu\,P^{(\pm)}$ with
$\mu\ge\|E_DNE_D\|\,\|(E_DPE_D)^{-1}\|$ works whenever $p$ and $1-p$ are boundedly
invertible on the relevant subspace, and in general one may take
$\tau_{++}=P\,\kappa\,P$, $\tau_{--}=-P^\perp\kappa P^\perp$ with $\kappa$ chosen so that
$E_D\tau_{++}E_D\ge\|E_DNE_D\|\cdot 1$ and $E_D\tau_{--}E_D\le-\|E_DNE_D\|\cdot1$.
Changing $\tau_{\pm\pm}$ changes neither $V(t)$ nor $\Tr(t\delta_c)$
(Lemma~\ref{lem:offdiag}).
\end{corollary}
\lby{Proposition~\ref{prop:split} and Lemma~\ref{lem:offdiag}; the displayed choices make
$\Sigma_2^{(0)}\ge0$ and $\Sigma_1^{(0)}\ge0$ by
$\pm E_DNE_D\le\|E_DNE_D\|\cdot1$}

\begin{verdictgap}[\;Corollary~\ref{cor:Pfree} is stated for bounded $E_DNE_D$]
The construction of $\tau_{\pm\pm}$ with $E_D\tau_{++}E_D\ge\lambda 1_D$ requires a
self-adjoint $\kappa$ supported in $I$ with $E_DP\kappa PE_D\ge\lambda 1_D$; since
$p=E_DPE_D$ has spectrum accumulating at $0$ and $1$ (the modular ultraviolet), such a
$\kappa$ must be unbounded, and then $t$ is not trace class and
Prop.~2.4 of \texttt{optimality\_all\_channels.tex} does not apply verbatim.  This is a
regularisation issue (cut the modular window at $|\kappa|\le\kappa_c$ and let
$\kappa_c\to\infty$), \emph{not} a structural obstruction; it is not pursued here
because condition (S) already decides the question.
\end{verdictgap}

\section{The isometric orbit and the obstruction}\label{sec:orbit}

The certificate can be tested without ever constructing $\Sigma_1,\Sigma_2$: it suffices
to look at one explicit two-parameter family of \emph{feasible} channels.

\begin{lemma}[The isometric orbit]\label{lem:orbit}
Let $\mathcal G$ be bounded, anti-self-adjoint, with $\mathcal G=E_D\mathcal GE_D$, and
$\eps\in\mathbb R$.  Put $\widetilde W_\eps:=W_s\,e^{-\eps\mathcal G}$.  Then
\begin{enumerate}[leftmargin=1.6em,itemsep=1pt]
\item $\widetilde W_\eps$ is a unitary of $L^2(\mathbb R)$, equal to the identity on
 $L^2(A)$, mapping $L^2(D)$ \emph{onto} $L^2(B)$; hence
 $V_\eps:=E_B\widetilde W_\eps E_D$ is unitary $\mathfrak h_D\to\mathfrak h_B$ and
 $(X,Z)=(V_\eps^*,V_\eps^*Q_{BB}V_\eps)\in\cF$;
\item the corresponding defect is
 $\delta(\eps)=E_I\bigl(P-e^{\eps\mathcal G}P_se^{-\eps\mathcal G}\bigr)E_I$,
 $P_s:=W_s^*PW_s$, and
 \begin{equation}\label{eq:deleps}
  \delta(\eps)=E_I\bigl[P,\ sD_\chi+\eps\mathcal G\bigr]E_I
   +O(s^2)+O(\eps^2)+O(\eps s);
 \end{equation}
\item consequently, with $u=\Pi M$ and $v_{\mathcal G}:=\Pi(P^\perp\mathcal GP)$,
 \begin{equation}\label{eq:gain}
  \min_{\eps}\ g_Q\bigl(\delta(\eps)\bigr)
   =\frac{s^2}{2}\Bigl(\|u\|_2^2-\frac{(\Re\langle u,v_{\mathcal G}\rangle)^2}{\|v_{\mathcal G}\|_2^2}\Bigr)+o(s^2)
   =g_Q(\delta_c)\bigl(1-\cos^2\angle(u,v_{\mathcal G})\bigr)+o(s^2).
 \end{equation}
\end{enumerate}
\end{lemma}
\begin{proof}
\lstep{1}{1}{(1).}
\lby{$e^{-\eps\mathcal G}$ is unitary, equals $1$ off $D$ and preserves $L^2(D)$
 (because $\mathcal G=E_D\mathcal GE_D$); $W_s$ is the push-forward along $k_s$ (C5),
 which is the identity on $A$ and maps $D$ onto $B$; a composition of such maps has the
 stated properties; feasibility as in (C5) since $V_\eps^*V_\eps=1$ makes both
 constraints active}
\lstep{1}{2}{(2).  $\widetilde W_\eps^*P\widetilde W_\eps=e^{\eps\mathcal G}P_se^{-\eps\mathcal G}$,
 and $Q_{\rm rec}=E_I\widetilde W_\eps^*P\widetilde W_\eps E_I$.}
\lby{$(W_se^{-\eps\mathcal G})^*=e^{\eps\mathcal G}W_s^*$; for the second claim compute the
 three blocks: $E_A\widetilde W^*P\widetilde WE_A=Q_{AA}$ (as $\widetilde WE_A=E_A$),
 $E_A\widetilde W^*P\widetilde WE_D=Q_{AB}V_\eps$, $E_D\widetilde W^*P\widetilde WE_D=Z$,
 which is \eqref{eq:F}}
\lstep{1}{3}{$P-e^{\eps\mathcal G}P_se^{-\eps\mathcal G}=(P-P_s)+\eps[P_s,\mathcal G]+O(\eps^2)$
 and $P-P_s=s[P,D_\chi]+O(s^2)$, $P_s=P+O(s)$.}
\lby{$\frac{\dd}{\dd\eps}e^{\eps\mathcal G}Ye^{-\eps\mathcal G}|_0=[\mathcal G,Y]$;
 (N5): $W_s=1-sD_\chi+O(s^2)$ gives $P_s=P+s[D_\chi,P]+O(s^2)$}
\lstep{1}{4}{(3).  For a positive quadratic form $g$ with bilinear form $b$,
 $\min_\eps g(\delta+\eps\eta)=g(\delta)-b(\delta,\eta)^2/g(\eta)$; here
 $g_Q(E_I[P,\mathcal Z]E_I)=\frac12\|\Pi P^\perp\mathcal ZP\|_2^2$ for anti-self-adjoint
 $\mathcal Z$, hence by polarisation
 $b\bigl(E_I[P,\mathcal Z_1]E_I,E_I[P,\mathcal Z_2]E_I\bigr)
  =\frac12\Re\langle\Pi P^\perp\mathcal Z_1P,\Pi P^\perp\mathcal Z_2P\rangle$.}
\lby{Prop.~\ref{prop:u}$\langle1\rangle2$ applied with $D_\chi$ replaced by
 $\mathcal Z$ (the derivation used only $\mathcal Z^*=-\mathcal Z$); the $O(\eps^2)$
 term of $\langle1\rangle3$ contributes $O(\eps^2 s)+O(\eps^4)=o(s^2)$ at the optimal
 $\eps=O(s)$}
\lqed{1}{5}
\end{proof}

\begin{theorem}[Obstruction $=$ non-self-adjointness of $E_DuE_D$]\label{thm:obstruction}
Let $R:=E_DuE_D$ and $R_a:=\frac12(R-R^*)$.  Then
\begin{equation}\label{eq:theta}
 \theta:=\sup_{\mathcal G}\cos^2\angle(u,v_{\mathcal G})
  =\frac{\langle R_a,\mathcal A^{-1}R_a\rangle_r}{\|u\|_2^2},\qquad
 \mathcal A(\mathcal G):=\tfrac12\bigl(E_Dv_{\mathcal G}E_D-(E_Dv_{\mathcal G}E_D)^*\bigr),
\end{equation}
$\langle\cdot,\cdot\rangle_r=\Re\Tr(\cdot^*\cdot)$, the supremum over bounded
anti-self-adjoint $\mathcal G=E_D\mathcal GE_D$.  In particular $\theta=0$ if and only if
$R_a=0$, i.e.\ if and only if the condition {\rm(S)} of Prop.~\ref{prop:split} holds; and
if $\theta>0$ then
\begin{equation}\label{eq:notopt}
 E^{(2)}\ \le\ (1-\theta)\,f_2z^2\,(1+o(1)),
\end{equation}
so the zero-collar compression is \emph{not} second-order optimal.
\end{theorem}
\begin{proof}
\lstep{1}{1}{$\Re\langle u,v_{\mathcal G}\rangle=\Re\langle u,P^\perp\mathcal GP\rangle
 =\Re\Tr(u^*\mathcal G)=\langle R,\mathcal G\rangle_r=\langle R_a,\mathcal G\rangle_r$.}
\lby{$u\in\overline{\mathcal N}$ and $v_{\mathcal G}=\Pi(P^\perp\mathcal GP)$ with $\Pi$ an
 orthogonal projection; $u=P^\perp uP$ so $\Tr(u^*P^\perp\mathcal GP)=\Tr(u^*\mathcal G)$;
 $\mathcal G=E_D\mathcal GE_D$ gives $\Tr(u^*\mathcal G)=\Tr((E_DuE_D)^*\mathcal G)$; and
 $\langle R_h,\mathcal G\rangle_r=0$ for $R_h$ self-adjoint and $\mathcal G$
 anti-self-adjoint}
\lstep{1}{2}{$\|v_{\mathcal G}\|_2^2=\langle\mathcal G,\mathcal A\mathcal G\rangle_r$, and
 $\mathcal A$ is a positive symmetric operator on the real space of such $\mathcal G$.}
\lby{$\|v\|^2=\Re\langle P^\perp\mathcal GP,v\rangle=\Re\Tr(\mathcal G^*E_DvE_D)
 =\langle\mathcal G,\mathcal A\mathcal G\rangle_r$ by the same manipulations as
 $\langle1\rangle1$; symmetry and positivity because
 $\langle\mathcal G_1,\mathcal A\mathcal G_2\rangle_r=\Re\langle v_1,v_2\rangle$}
\lstep{1}{3}{\eqref{eq:theta} and ``$\theta=0\iff R_a=0$''.}
\lby{$\sup_x\langle b,x\rangle^2/\langle x,\mathcal Ax\rangle=\langle b,\mathcal A^{-1}b\rangle$
 (on $\overline{\operatorname{ran}}\mathcal A$; $=+\infty$ if $b\notin\overline{\operatorname{ran}}\mathcal A$,
 which cannot happen since $\theta\le1$ by Cauchy--Schwarz); $\theta=0$ iff the linear
 functional $\langle R_a,\cdot\rangle_r$ vanishes, i.e.\ iff $R_a=0$ (take
 $\mathcal G=R_a$)}
\lstep{1}{4}{\eqref{eq:notopt}.}
\lby{Lemma~\ref{lem:orbit}(3), the feasibility in Lemma~\ref{lem:orbit}(1), and
 Prop.~\ref{prop:u}(3)}
\lqed{1}{5}
\end{proof}

\section{Numerics: $\theta$ is not zero}\label{sec:num}

\paragraph{Model.}  $L^2(S^1)$ on an $L$-point grid in the NS sector (half-integer Fourier
modes $\nu=n+\frac12$), $P=\mathbf 1_{\nu>0}$ an \emph{exact} projection, $E_S$ exact
diagonal projections.  Cayley chart $x=\tan(\theta/2)$: the line configuration
$A=(-a,0)$, $D=(0,L_g)$, $I'=\mathbb R\setminus\overline I$ becomes
$\theta_A=-2\arctan a$, $\theta_D=2\arctan L_g$, and the line M\"obius field $-x^2\partial_x$
becomes $\chi(\theta)=\cos\theta-1$; $\chi=0$ on $A$ and
$\chi=(\cos\theta-1)e^{-\alpha(\theta-\theta_D)^2}$ on $I'$ (a $C^\infty$ admissible
extension, $\alpha$ free).  $D_\chi=\frac12(\chi\partial+\partial\chi)$ with a
\emph{smoothstep ultraviolet taper} on the spectral derivative, which switches off the
spurious second Fermi point of the periodic grid.  Scripts:
\texttt{numerics/optimality\_all/s10\_circle2.py} (validation),
\texttt{s10\_theta.py}, \texttt{s10\_diag.py} ($\theta$ by conjugate gradients on
\eqref{eq:theta}), \texttt{s10\_nonpert.py} (non-perturbative check, output
\texttt{s10\_nonpert.out}).

\paragraph{(a) Validation of the framework.}  With $u$ obtained from
$\mathcal S_Q(\ell)=-\dot\delta$, $u=P^\perp\ell P$:
$\|u\|_2^2$ is \emph{independent of the extension} $\alpha\in\{1.5,3,6\}$ to $4$ digits
and of the weight cap over $10^{4}$--$10^{12}$, whereas $\|M\|_2^2$ varies by $15\%$ --
exactly as Prop.~\ref{prop:u}(1) predicts.  Quantitatively, with
$\zeta=aL_gs/(a+L_g)$ and $f_2=1/(12\pi^2)$,
\[
 \frac{\tfrac12 s^2\|u\|_2^2}{f_2\zeta^2}=0.984,\ 0.987,\ 0.989,\ 0.991
 \quad\text{for }L=256,384,512,768\ (a=2,L_g=1),
\]
and $0.984$--$0.992$ for $(a,L_g)=(1,1),(4,1),(8,1),(2,\tfrac12)$: the machinery
reproduces Theorem~A (\texttt{quadratic\_limit.tex} Thm.~4.2 with
\texttt{uhlmann\_upper\_bound.tex} Thms.~8.1--8.2) to $1\%$, with a clean $O(1/L)$ drift.

\paragraph{(b) The obstruction.}
$R=E_DuE_D$ is, to the accuracy reached, \emph{purely anti-self-adjoint}:
$\|R_a\|_2/\|R\|_2=0.976,\,0.989,\,0.994,\,0.997$ for $L=256,384,512,768$.  Condition
{\rm(S)} of Prop.~\ref{prop:split} therefore fails maximally.

\paragraph{(c) The gain.}  Conjugate gradients on \eqref{eq:theta} (converged: the CG
ladder is flat after $\sim60$ iterations) gives
\begin{center}\small
\begin{tabular}{lccccc}
\toprule
 & $L=256$ & $L=384$ & $L=512$ & $(a,L_g)=(1,1)$ & $(4,1)$\\
\midrule
$\theta$ & $0.2895$ & $0.2924$ & $0.2950$ & $0.2999$ & $0.2918$\\
\bottomrule
\end{tabular}
\end{center}
Richardson extrapolation in $1/L$ gives $\theta_\infty=0.300\pm0.005$.  The value is
independent of the geometry to within the discretisation error, as it must be: two
adjacent intervals on the circle are a single M\"obius orbit.  It is stable under the
ultraviolet taper ($\theta=0.2894$ and $0.2946$ for tapers starting at $0.20L$ and
$0.30L$; the aggressive taper at $0.38L$ is contaminated -- it also spoils
$\|u\|_2^2$ by $4\%$ -- and is discarded) and under the weight cap
($10^{6}$ vs $10^{12}$: $0.29462$ vs $0.29459$).  The optimiser is a \emph{bounded}
rotation: $\|\eps_*\mathcal G_*\|_{\rm op}/s=0.49$, $\|\eps_*\mathcal G_*\|_2/s=0.84$.

\paragraph{(d) Non-perturbative check.}  Building $W_s=e^{-sD_\chi}$ (exactly unitary),
$\widetilde W=W_se^{-\eps\mathcal G_*}$, the \emph{exact} defect
$\delta(\eps)=E_I(P-\widetilde W^*P\widetilde W)E_I$ and the \emph{exact} quadratic form
$g_Q$ (cap $10^4$, $L=512$, $a=2$, $L_g=1$):
\begin{center}\small
\begin{tabular}{lccc}
\toprule
$s$ & $g_Q(\delta_c)\big/\tfrac12s^2\|u\|_2^2$ & $\min_\eps g_Q(\delta(\eps))/g_Q(\delta_c)$ & optimal $\eps/s$\\
\midrule
$0.001$ & $0.9985$ & $0.7050$ & $-1.000$\\
$0.002$ & $0.9983$ & $0.7046$ & $-1.000$\\
$0.005$ & $0.9978$ & $0.7038$ & $-1.000$\\
\bottomrule
\end{tabular}
\end{center}
against the prediction $1-\theta=0.7054$: agreement to $10^{-3}$, with the minimiser at
the predicted location.  (At $s\ge0.02$ with an uncapped weight the comparison is
destroyed by the order-of-limits effect of S8's Verdict box -- $g_Q(\delta_c)$ then
exceeds $\tfrac12s^2\|u\|^2$ by factors $2.8$, $12$, $45$ at $s=0.02,0.05,0.1$ -- which
is why the check must be run at $s\lesssim5\times10^{-3}$.)

\section{What the obstruction is, structurally}\label{sec:struct}

(Section~\ref{sec:exact} proves the exact identity behind the numerical observation of
\S\ref{sec:num}(b); this section explains the mechanism first.)

\begin{lemma}[The model computation]\label{lem:model}
Let $\mathcal G$ be bounded anti-self-adjoint with $\mathcal G=E_D\mathcal GE_D$, and
$p:=E_DPE_D$.  Then, with $N_{\mathcal G}:=P^\perp\mathcal GP$,
\begin{equation}\label{eq:model}
 E_DN_{\mathcal G}E_D-\bigl(E_DN_{\mathcal G}E_D\bigr)^*
 = p\,\mathcal G\,(1-p)+(1-p)\,\mathcal G\,p=\mathcal S_p(\mathcal G),
\end{equation}
which vanishes if and only if $\mathcal G=0$ (because $0<p<1$).
\end{lemma}
\lby{$\mathcal G=E_D\mathcal GE_D$ gives $E_DP^\perp\mathcal GPE_D=(E_DP^\perp E_D)\mathcal G(E_DPE_D)=(1-p)\mathcal Gp$
and $(E_DP^\perp\mathcal GPE_D)^*=E_DP\mathcal G^*P^\perp E_D=-p\,\mathcal G\,(1-p)$; the
spectrum of the modular operator of an interval fills $(0,1)$, so
$\mathcal S_p$ is injective}

\begin{remark}[Reading of Lemma~\ref{lem:model} and of the numerics]\label{rem:reading}
Lemma~\ref{lem:model} is the ``local'' version of the obstruction: for a rotation
generator \emph{strictly inside} $D$ the $D$-corner of the off-diagonal block is never
self-adjoint.  The compression direction is not of that form -- $D_\chi$ is a
differential operator whose vector field does \emph{not} vanish at the right endpoint of
$D$ (that is exactly why the compression shrinks $I$ and a unitary of $\mathfrak h_D$
cannot imitate it), and the relevant object is the projected $u=\Pi M$, not $M$.  The
numerics of \S\ref{sec:num} say that the projection does \emph{not} restore
self-adjointness: $E_DuE_D$ is, in the continuum limit, purely anti-self-adjoint.
Equivalently, and this is the sharpest way to say it: \emph{the compression direction
$u$ has an overlap $\cos^2\angle=\theta\approx0.30$ with the space of directions
generated by Bogoliubov rotations inside $B$, so $30\%$ of the compression defect can be
rotated away.}
\end{remark}

\section{Consequences for the dual value and for the recovery error}\label{sec:conseq}

\begin{theorem}[Second-order value of the programme]\label{thm:value}
Assume the numerical value $\theta=0.300\pm0.005>0$ of \S\ref{sec:num}.  Then, for the
chiral free fermion,
\begin{equation}\label{eq:final}
 E^{(2)}=\min_{(X,Z)\in\cF}g_Q\bigl(Q-Q_{\rm rec}(X,Z)\bigr)\ \le\ (1-\theta)\,f_2z^2(1+o(1))
 \ =\ (0.70\pm0.01)\,f_2z^2(1+o(1)),
\end{equation}
and consequently
\begin{enumerate}[leftmargin=1.6em,itemsep=1pt]
\item the zero-collar geometric compression is \textbf{not} the second-order optimum of
 the quasi-free programme; the KKT certificate of Theorem~\ref{thm:kkt} does \emph{not}
 exist at the compression, because condition {\rm(S)} fails;
\item the Sion dual value satisfies $\mathcal D_z\le E^{(2)}\le0.71f_2z^2$, so
 \emph{no} test operator $t$ can certify $\Erecc\ge f_2z^2$: route~2 of the Phase-3b brief
 is closed, and so is the programme
 ``\,$\liminf z^{-2}\Erecc\ge f_2$\,'' by way of the quasi-free relaxation;
\item if in addition the optimal rotation $e^{\eps\mathcal G_*}$ is implementable in the
 Fock representation (Shale--Stinespring: $[P,\mathcal G_*]\in\Stwo$), then the channel
 $\beta_*=\operatorname{Ad}\Gamma(W_s e^{-\eps\mathcal G_*})|_{\cA(D)}$ is an admissible
 recovery channel and $\Erecc\le(1-\theta)f_2z^2(1+o(1))$: the geometric compression is
 beaten by an actual channel, and $\kappa_{\rm opt}\le0.71$ in the notation of S8's
 Verdict~Q3.
\end{enumerate}
\end{theorem}
\begin{proof}
\lstep{1}{1}{\eqref{eq:final}.}
\lby{Lemma~\ref{lem:orbit}(1) exhibits feasible points of $\cF$;
 Lemma~\ref{lem:orbit}(3) and Theorem~\ref{thm:obstruction} evaluate the objective there;
 Prop.~\ref{prop:u}(3) identifies $g_Q(\delta_c)=f_2z^2(1+o(1))$}
\lstep{1}{2}{(1).}
\lby{if the certificate \eqref{eq:cert} held for $t=t_c$ then, by
 Theorem~\ref{thm:kkt} and Corollary~\ref{cor:test}, $(X_c,Z_c)$ would minimise, so
 $E^{(2)}=g_Q(\delta_c)=f_2z^2(1+o(1))$, contradicting \eqref{eq:final}}
\lstep{1}{3}{(2).}
\lby{\texttt{optimality\_all\_channels.tex} Thm.~4.1$\langle1\rangle3$:
 $\mathcal D_z\le E^{(2)}$ for every $t$; and \eqref{eq:final}}
\lstep{1}{4}{(3).}
\lby{Shale--Stinespring (Araki 1987, Thm.~6.13; Baumg\"artel--Jurke--Lled\'o 2002,
 Thm.~4.13) makes $\alpha_{e^{-\eps\mathcal G_*}}$ unitarily implementable in the Fock
 representation; $W_s$ is implementable by
 \texttt{uhlmann\_upper\_bound.tex} Prop.~1.7; the composite $\widetilde W$ is then
 implementable, $\operatorname{Ad}\Gamma(\widetilde W)$ restricts to a normal unital
 $*$-endomorphism $\cA(D)\to\cA(B)$ (Lemma~\ref{lem:orbit}(1)), hence to an admissible
 recovery channel in the sense of S8 Def.~1.1; its recovered symbol is
 $Q_{\rm rec}=E_I\widetilde W^*P\widetilde WE_I$ by Lemma~\ref{lem:orbit}(2) and its
 error is $g_Q(\delta(\eps))(1+o(1))$ by Note~5 Prop.~2.3}
\lqed{1}{5}
\end{proof}

\paragraph{(e) Implementability of the optimal rotation.}
\texttt{s10\_implement.py} gives, for $L=256,384,512$,
\[
 \|P^\perp(\eps_*\mathcal G_*)P\|_2/s=0.0608,\ 0.0625,\ 0.0638,\qquad
 \|\eps_*\mathcal G_*\|_{\rm op}/s=0.482,\ 0.472,\ 0.486,
\]
both bounded and converging: $[P,\mathcal G_*]\in\Stwo$, so the Shale--Stinespring
condition of Theorem~\ref{thm:value}(3) is met.  The single explicit direction
$\mathcal G=R_a$ already gives $\cos^2\angle(u,v_{R_a})=0.1234,\ 0.1242,\ 0.1250$
(stable): even the crudest choice beats the compression by $12\%$.

\section{Route 3: what a rigorous lower bound can still give}\label{sec:lower}

By Theorem~\ref{thm:value}(2) the best any dual test operator can certify is
$\mathcal D_z\le E^{(2)}\le0.71f_2z^2$, so the target $f_2z^2$ is unreachable and the
correct target is $E^{(2)}$ itself.  Two remarks.

\begin{proposition}[The isometric orbit computes its own value exactly]\label{prop:orbval}
Let $\mathcal Y:=\overline{\{\Pi(P^\perp\mathcal GP):\mathcal G^*=-\mathcal G\ \text{bounded},\ \mathcal G=E_D\mathcal GE_D\}}$.
The minimum of $g_Q$ over the isometric orbit
$\{(V^*,V^*Q_{BB}V):V=E_B\widetilde WE_D,\ \widetilde W\ \text{as in Lemma~\ref{lem:orbit}(1)}\}$
equals $\frac{s^2}{2}\operatorname{dist}(u,\mathcal Y)^2+o(s^2)=(1-\theta)f_2z^2(1+o(1))$,
and the minimiser is unique up to the kernel of $\Pi$.
\end{proposition}
\lby{Lemma~\ref{lem:orbit}(2)--(3): on the orbit
$\delta=E_I[P,sD_\chi+\mathcal G]E_I+o(s)$ with $\mathcal G$ arbitrary in the class, and
$g_Q(E_I[P,\mathcal Z]E_I)=\frac12\|\Pi P^\perp\mathcal ZP\|_2^2$
(Prop.~\ref{prop:u}$\langle1\rangle2$), so the minimisation is the orthogonal projection
of $su$ onto $\mathcal Y$; $\theta=\|\Pi_{\mathcal Y}u\|^2/\|u\|^2$ by
Theorem~\ref{thm:obstruction}}

\begin{verdictgap}[\;The non-isometric directions were not analysed]
$\cF$ also contains points with $XX^*<1$ and genuine noise $0<Y\le1-XX^*$, i.e.\
channels that are not $*$-endomorphisms.  Those were not included: $(1-\theta)f_2z^2$ is
an \emph{upper} bound for $E^{(2)}$, not its value.  A matching lower bound would
require the full first-order analysis of $\cF$ at the new optimum (whose location is not
known in closed form).  In S8's lattice numerics the optimal $X$ has singular values
$0.95$--$1.00$, i.e.\ the optimum is nearly isometric, which suggests that the isometric
orbit carries most of the gain, but this is not proved.
\end{verdictgap}

\begin{verdictgap}[\;From $g_Q$ back to $-\log\Fid$ for the rotated channel]
Theorem~\ref{thm:value}(3) asserts $\Erecc\le(1-\theta)f_2z^2(1+o(1))$.  This uses
$-\log\Fid=g_Q(\delta)(1+o(1))$ for the rotated channel, i.e.\ the \emph{upper} half of
the second-order expansion.  In the continuum that half is exactly what
\texttt{uhlmann\_upper\_bound.tex} had to prove separately for the compression
(Thms.~8.1--8.2), and it is not automatic (S8's ``order of limits'' Verdict).  The
lower half, $-\log\Fid\ge g_Q(\delta)(1-o(1))$, holds for every channel
(\texttt{optimality\_all\_channels.tex} Prop.~2.4), so the \emph{programme} statement
\eqref{eq:final} and the blocking statement Theorem~\ref{thm:value}(2) are unaffected;
only the claim that an actual channel realises $0.70f_2z^2$ needs this extra step.
\end{verdictgap}

\begin{remark}[Tension with S8's lattice numerics, and how to resolve it]\label{rem:tension}
S8's Verdict~Q3 reports $\kappa_{\rm opt}=1.00\pm0.15$ from six lattice geometries.
That is not in contradiction with $\kappa_{\rm opt}\le0.71$: S8's channels are produced
by minimising a \emph{capped} $g_Q$ ($\kappa_c\le12$--$20$) and are then evaluated by the
exact $-\log\Fid$; the cap removes precisely the modular ultraviolet at the endpoints of
$D$ where $\mathcal G_*$ lives (the optimal rotation is bounded, but its Riesz vector
$v_{\mathcal G_*}$ is spread over the whole modular window).  Moreover at $z\simeq1/14$
the unrestricted $g_Q$ overshoots $-\log\Fid$ by a factor $\simeq28$, so the search is
not minimising the quantity being reported.  The decisive test is S11's: take the lattice
$(X_c,Z_c)$ closest to the compression, form $t_c$, and check whether
$\Tr(t_cQ_{\rm rec})$ is maximised there --- Theorem~\ref{thm:kkt} says it is not, and
predicts the maximiser to be an (almost) isometric $X_cU$.
\end{remark}

\section{An exact identity: $E_DuE_D$ is purely anti-self-adjoint}\label{sec:exact}

The numerical observation of \S\ref{sec:num}(b) has an exact explanation, which upgrades
the obstruction from ``numerically nonzero'' to ``analytically maximal''.

\begin{lemma}[M\"obius rigidity of the $DD$-block]\label{lem:rigid}
As a sesquilinear form on $C_c^\infty(D)$,
\begin{equation}\label{eq:rigid}
 E_D\,\dot\delta\,E_D=E_D[P,D_\chi]E_D=[\,p\,,\,d\,]=0,\qquad p:=E_DPE_D,\quad d:=E_DD_\chi E_D ,
\end{equation}
i.e.\ \emph{the $DD$-block of the compression defect vanishes at first order in $s$}.
\end{lemma}
\begin{proof}
\lstep{1}{1}{\lpick{} $f,g\in C_c^\infty(D)$.  Then $E_Dg=g$, $D_\chi g$ is supported in
 $\operatorname{supp}g\subset D$, and $\chi|_D$ coincides with the \emph{global} M\"obius
 field $\chi_{\rm glob}(x)=-x^2$ (in the normalisation of (N5)); hence
 $D_\chi g=D_{\rm glob}g$ and $E_DD_\chi g=D_{\rm glob}g$.}
\lby{$D_\chi$ is a first-order differential operator, hence local, and $\chi=\chi_{\rm glob}$
 on $D$ by (N5)}
\lstep{1}{2}{$\langle f,p\,d\,g\rangle=\langle f,PD_{\rm glob}g\rangle$ and
 $\langle f,d\,p\,g\rangle=\langle f,D_{\rm glob}Pg\rangle$.}
\lby{$\langle1\rangle1$ and $E_Df=f$ for the first; for the second,
 $\langle f,E_DD_\chi E_DPE_Dg\rangle=\langle f,D_\chi(E_DPg)\rangle
 =-\langle D_\chi f,E_DPg\rangle=-\langle D_\chi f,Pg\rangle=\langle f,D_{\rm glob}Pg\rangle$,
 using $D_\chi^*=-D_\chi$, that $f$ has compact support inside $D$ (so no boundary term
 and the distributional jumps of $E_DPg$ at $\partial D$ are annihilated), and
 $D_\chi f=D_{\rm glob}f$ supported in $D$}
\lstep{1}{3}{$[P,D_{\rm glob}]=0$.}
\lby{$D_{\rm glob}=\chi_{\rm glob}\partial_x+\frac12\chi_{\rm glob}'$ with
 $\chi_{\rm glob}(x)=-x^2$ is the generator of the one-parameter M\"obius group
 $x\mapsto x/(1+sx)$ acting on $L^2(\mathbb R)$ by the half-density push-forward
 \eqref{eq:pf}; that action preserves the Hardy space, i.e.\
 $W_{h_s}^*PW_{h_s}=P$ for all $s$ (M\"obius covariance of the chiral vacuum,
 \texttt{uhlmann\_upper\_bound.tex} \S6 ``the circle chart''; equivalently the Cayley
 transform sends $D_{\rm glob}$ to $L_{1}-L_{-1}$-type generators that annihilate the
 vacuum representation's positive-frequency subspace).  Differentiating at $s=0$ gives
 $[D_{\rm glob},P]=0$}
\lstep{1}{4}{\eqref{eq:rigid}.}
\lby{$\langle1\rangle2$--$\langle1\rangle3$; $\dot\delta$ is Hilbert--Schmidt
 (Prop.~\ref{prop:u}), so a bounded operator, and a bounded operator whose form vanishes
 on the dense set $C_c^\infty(D)$ has $E_D\dot\delta E_D=0$}
\lqed{1}{5}
\end{proof}

\begin{corollary}[The obstruction is maximal]\label{cor:maximal}
$E_DuE_D$ is purely anti-self-adjoint:
$E_DuE_D+(E_DuE_D)^*=E_D\dot\delta E_D\cdot(-1)=0$.  Hence condition {\rm(S)} of
Prop.~\ref{prop:split} holds if and only if $E_DuE_D=0$, and
\begin{equation}\label{eq:thetalb}
 \theta\ \ge\ \frac{\|E_DuE_D\|_2^2}{\|u\|_2^2}\ (>0\ \text{numerically}:\ 0.0143).
\end{equation}
\end{corollary}
\begin{proof}
\lstep{1}{1}{$E_I(u+u^*)E_I=\mathcal S_Q(\ell_*)=-\dot\delta$, where $u=P^\perp\ell_*P$ and
 $\mathcal S_Q(\ell_*)=-\dot\delta$ is the normal equation defining $u=\Pi M$.}
\lby{$\mathcal S_Q(\ell)=E_I(P\ell P^\perp+P^\perp\ell P)E_I$ for $\ell=E_I\ell E_I$ (N4),
 and $P\ell P^\perp=(P^\perp\ell P)^*=u^*$; the normal equation is
 $\Re\langle N_\ell,\Pi M\rangle=\Re\langle N_\ell,M\rangle$ for all $\ell$, i.e.\
 $\frac12\Tr(\ell\,\mathcal S_Q(\ell_*))=-\frac12\Tr(\ell\dot\delta)$ by
 Lemma~\ref{lem:offdiag}$\langle1\rangle2$ and Prop.~\ref{prop:u}$\langle1\rangle1$}
\lstep{1}{2}{Compress $\langle1\rangle1$ with $E_D$ and use Lemma~\ref{lem:rigid}.}
\lby{$E_DE_I=E_D$}
\lstep{1}{3}{\eqref{eq:thetalb}.}
\lby{Theorem~\ref{thm:obstruction} with the test direction $\mathcal G=R_a=E_DuE_D$,
 and $\|v_{\mathcal G}\|\le\|P^\perp\mathcal GP\|\le\|\mathcal G\|_2$}
\lqed{1}{4}
\end{proof}

\begin{verdictbreak}[\;(S) fails; the compression is not the optimum of the programme]
Lemma~\ref{lem:rigid} and Corollary~\ref{cor:maximal} reduce the certificate to the
single question ``$E_DuE_D=0$?''.  The answer is no: $\|E_DuE_D\|_2^2/\|u\|_2^2=0.0143$
in the spectral model of \S\ref{sec:num}, stable in $L$, in the weight cap, in the
ultraviolet taper and in the geometry, and the full supremum is
$\theta=0.300\pm0.005$.  The non-perturbative check \S\ref{sec:num}(d) reproduces
$1-\theta$ to $10^{-3}$ from the exact defect of an explicitly constructed feasible
channel.  Consequently the KKT certificate of Theorem~\ref{thm:kkt} does not exist at the
zero-collar compression, and \eqref{eq:final} holds.  The one ingredient that is
numerical rather than proved is the \emph{non-vanishing} of $E_DuE_D$ (equivalently
$\theta>0$); everything else in the chain is proved above.
\end{verdictbreak}

\section{Verdict}\label{sec:verdict}

\paragraph{Proved here (unconditionally).}
\begin{enumerate}[leftmargin=1.8em,itemsep=1pt]
\item[\textbf{P1}] \emph{Exact certificate.} Theorem~\ref{thm:kkt}: in the continuum,
 where the zero-collar compression is a \emph{unitary} $V_c:\mathfrak h_D\to\mathfrak h_B$,
 the point $(X_c,Z_c)$ maximises $\Tr(t\,Q_{\rm rec})$ over $\cF$ \emph{if and only if}
 the uniquely determined multiplier
 $\Sigma_1=t_{DA}Q_{AB}V_c-t_{DD}V_c^*R_-V_c$ is self-adjoint and
 $\Sigma_1\ge0$, $\Sigma_1+t_{DD}\ge0$.  Necessity uses only explicit feasible
 perturbations (no constraint qualification); sufficiency is weak duality for a concave
 Lagrangian.  Both LMI constraints are active at the compression and the multipliers are
 unique.
\item[\textbf{P2}] \emph{Legendre reduction.} Lemma~\ref{lem:offdiag}: for a test
 operator supported in $I$, $V(t)=\|P^\perp tP\|_2^2$ and $\Tr(t\delta_c)$ depend only on
 $N=P^\perp tP$; the two $P$-diagonal blocks are free.  Prop.~\ref{prop:split}: the
 certificate then splits into (S) ``$E_DNE_D$ self-adjoint'' -- which the free blocks
 cannot influence -- and (P) a positivity sandwich which they can always satisfy.
\item[\textbf{P3}] \emph{Value of the compression, intrinsically.}
 Prop.~\ref{prop:u}: $g_Q(\delta_c)=\frac{s^2}2\|\Pi M\|_2^2+O(s^3)$ with
 $\Pi M=u$ independent of the $C^{1,1}$ extension, and
 $\|u\|_2^2=g_B/4$, i.e.\ $g_Q(\delta_c)=f_2z^2(1+o(1))$ -- Theorem~A recovered in the
 present variables and confirmed numerically to $1\%$.
\item[\textbf{P4}] \emph{The isometric orbit.} Lemma~\ref{lem:orbit}: for every bounded
 anti-self-adjoint $\mathcal G$ supported in $D$, $\widetilde W=W_se^{-\eps\mathcal G}$
 gives a feasible point of $\cF$ with
 $\min_\eps g_Q=g_Q(\delta_c)(1-\cos^2\angle(u,v_{\mathcal G}))$, and
 Theorem~\ref{thm:obstruction}: $\sup_{\mathcal G}\cos^2\angle=\theta$ vanishes iff
 $E_DuE_D$ is self-adjoint, i.e.\ iff (S) holds.
\item[\textbf{P5}] \emph{M\"obius rigidity.} Lemma~\ref{lem:rigid}: the $DD$-block of the
 compression defect vanishes at first order, $E_D\dot\delta E_D=[p,d]=0$, because the
 M\"obius generator commutes with the Hardy projection and is local.  Hence
 (Cor.~\ref{cor:maximal}) $E_DuE_D$ is \emph{purely anti-self-adjoint}: condition (S) is
 equivalent to $E_DuE_D=0$, so the certificate either exists trivially or fails
 maximally.
\end{enumerate}

\paragraph{Established numerically (spectral model, all diagnostics stable).}
\begin{enumerate}[leftmargin=1.8em,itemsep=1pt]
\item[\textbf{N1}] $E_DuE_D\neq0$: $\|E_DuE_D\|_2^2/\|u\|_2^2=0.0143$.
\item[\textbf{N2}] $\theta=0.300\pm0.005$, geometry independent (as M\"obius invariance
 demands), cap independent, extension independent, ultraviolet-taper independent.
\item[\textbf{N3}] Non-perturbative confirmation: the exact defect of the explicitly
 constructed feasible channel $\widetilde W=W_se^{-\eps\mathcal G_*}$ has
 $\min_\eps g_Q/g_Q(\delta_c)=0.705$ at $s=10^{-3}$--$5\times10^{-3}$, matching
 $1-\theta$ to $10^{-3}$.
\item[\textbf{N4}] $[P,\mathcal G_*]\in\Stwo$ (Shale--Stinespring), so the optimal
 rotation is implementable and $\beta_*$ is an honest recovery channel.
\end{enumerate}

\begin{verdictbreak}[\;Answer to the S10 question: \textbf{No}]
The zero-collar geometric compression is \emph{not} the second-order optimum of the
quasi-free convex programme.  Modulo the single numerical input N1
($E_DuE_D\neq0$, equivalently $\theta>0$),
\[
 E^{(2)}=\min_{(X,Z)\in\cF}g_Q\bigl(Q-Q_{\rm rec}(X,Z)\bigr)\ \le\ (1-\theta)f_2z^2(1+o(1))
 \ =\ (0.70\pm0.01)\,f_2z^2(1+o(1)) ,
\]
attained (to that order) by \emph{composing the compression with a Bogoliubov rotation
$e^{\eps\mathcal G_*}$ inside $B$}, $\eps=O(s)$, $\|\eps\mathcal G_*\|_{\rm op}=O(s)$.
The mechanism is transparent: the compression defect direction $u$ has an angle with the
span of the rotation directions whose cosine squared is $\theta\approx0.3$, so $30\%$ of
the second-order defect can be rotated away inside $B$; M\"obius rigidity
(Lemma~\ref{lem:rigid}) guarantees that the part of $E_DuE_D$ that would have to vanish
for optimality is precisely the part that survives.
\end{verdictbreak}

\begin{verdictbreak}[\;Consequence for the Phase-3b programme: route 2 is closed]
By \texttt{optimality\_all\_channels.tex} Thm.~4.1, $\mathcal D_z\le E^{(2)}$, so
\emph{no} test operator $t$ -- the compression's SLD included -- can certify
$\liminf z^{-2}\Erecc\ge f_2$.  The target of the Phase-3b brief
(``the programme has second-order value $f_2z^2$'') is false as stated; the correct
target is $E^{(2)}$, which is at most $0.71f_2z^2$.  The all-channel equivalence of
Thm.~4.1(3) is untouched: it says the all-channel and quasi-free second-order optima
coincide, and both are now known to be strictly below the compression value.
\end{verdictbreak}

\begin{verdictgap}[\;What is not settled]
(i) $\theta>0$ is numerical, not proved; a proof needs
$E_DuE_D\ne0$, i.e.\ that the Riesz vector of the compression direction has a nonzero
$DD$-corner.  (ii) $(1-\theta)f_2z^2$ is an \emph{upper} bound for $E^{(2)}$: the
non-isometric part of $\cF$ (genuine noise $Y>0$) was not analysed, so
$\kappa_{\rm opt}\le0.71$ but its exact value is open.  (iii) The passage from
$g_Q$ to $-\log\Fid$ for the rotated channel (the ``upper half'' of the second-order
expansion) is assumed, as discussed in \S\ref{sec:lower}; the programme statement and
the blocking statement do not need it.  (iv) Theorem~\ref{thm:kkt} was stated for
trace-class $t$; the regularisation of the diagonal blocks in Cor.~\ref{cor:Pfree} was
not carried out (it is not needed for the refutation).
\end{verdictgap}
\end{document}
